\subsection{The first nonzero class quotient after two inert support primes}

\begin{statement}{A.151}{Proved}
{An explicit compatible ray-class quotient in the two-prime tower}
Let \(M_1,M_2\) be positive multiples of twenty-four,
and put
\[
 \begin{gathered}
 k=\Q(\mu_3),\qquad d=11\cdot59=649,
 \qquad F=k(d^{1/3}),\\
 a=3\cdot11^{M_1}59^{M_2},\qquad
 L_n=\Q(\mu_{3^n},a^{1/3^n}),\\
 n_0=1+\max\{v_3(M_1),v_3(M_2)\}.
 \end{gathered}
\tag{A.151.1}
\]
The roots defining the \(L_n\) are compatible.
Let \(S_n\) consist of their places above
\(2,3,11,59,\infty\), and set
\(\mathcal C_n=\operatorname{Cl}(\mathcal O_{L_n,S_n})\{3\}\).
For every \(n\geq n_0\), \(FL_n/L_n\) is a
nontrivial cyclic cubic extension, unramified at
every finite place and split at every place in
\(S_n\).
Its Artin character gives compatible surjections
\[
 \chi_n:\mathcal C_n\longrightarrow
               \operatorname{Gal}(F/k)=\Z/3,
 \qquad
 \chi_n\circ N_{L_{n+1}/L_n}=\chi_{n+1}.
\tag{A.151.2}
\]
The induced map
\[
                 \varprojlim_{n\geq n_0}\mathcal C_n
                          \longrightarrow\Z/3
\tag{A.151.3}
\]
is onto.
In particular, the class-group vanishing mechanism
of A.149 does not persist in these specified
two-prime Galois layers.

The example retains the finite old-cover residue
condition: \(-2\) is a twenty-fourth power modulo
both eleven and fifty-nine.  When \(M_1=M_2=24\),
the cover \(aX^{24}-2Y^{24}=1\) has unit-coordinate
points over every \(\Q_p\).
No value of a particular target norm obstruction is
deduced from (A.151.3).  The theorem gives the actual
compatible class character against which such a
target can next be paired.
\end{statement}
\begin{proof}
\emph{The precise two-prime ray field.}
Both eleven and fifty-nine are inert in
\(k=\Q(\mu_3)\), and
\[
 v_3(11^2-1)=v_3(59^2-1)=1.
\tag{A.151.4}
\]
The two residue-unit three-parts are each \(\mu_3\).
Since \(\operatorname{Cl}(k)=0\) and
\(\mathcal O_k^\times=\mu_6\), the ray exact
sequence of \cite[V, Theorem 1.7]{MilneCFT}
identifies the three-primary ray group at
\((11)(59)\) with
\[
                   (\Z/3\oplus\Z/3)/\langle(1,1)\rangle
                                      \simeq\Z/3.
\tag{A.151.5}
\]
The field \(F=k(649^{1/3})\) is cyclic of degree
three, since its radicand has valuation one at
both of these primes.
It is tamely ramified at both, so their conductor
exponents are one.
It is unramified away from \(3,11,59\) by its
Kummer equation.
At three, \(649\equiv1\pmod9\), and more
explicitly
\(v_3(649-1)=4>2v_3(3)\).
Hensel's lemma applied at one makes it an actual
cube in \(\Q_3\).
Thus every prime over three splits in \(F/k\).
The same is true over two: the cube map is
invertible on \(\Z_2^\times\).
Consequently \(F\) is exactly the nontrivial
three-primary ray field of (A.151.5), by
\cite[V, Theorems 3.5--3.6]{MilneCFT}.

The split condition at three is also visible
in the ray map itself.
For \(\zeta\in\mu_3\) and \(\pi=1-\zeta\),
\(\pi^2=-3\zeta\).
Every rational residue unit is a cube at
either inert prime, because both rational
residue groups have order prime to three.
Normalize each residue cubic character by
\(\lambda_i(\zeta)=1\).
Then \(2\lambda_i(\pi)=1\), so both
\(\lambda_i(\pi)=2\).
The ray character \((x,y)\mapsto x-y\)
therefore kills the three-prime, as well
as the principal ideal \((2)\).
This is the explicit class-field character
used below.

\emph{The global intersection through the three-adic degree.}
The rational radicand \(a\) has three-adic
valuation one.
Over \(\Q_3(\mu_{3^n})\) it is not a cube:
\(\Q_3(a^{1/3})/\Q_3\) is Eisenstein of
degree three and not Galois, since it cannot
contain the degree-two field \(\Q_3(\mu_3)\).
It cannot be a subfield of an abelian local
cyclotomic extension.
Kummer theory therefore gives the full
\(3^n\)-degree radical extension over that
cyclotomic field.
The local and global closures both have
the actual semidirect group
\[
              C_{3^n}\rtimes(\Z/3^n)^\times.
\tag{A.151.6}
\]
The complement fixes the chosen radical,
and its element with cyclotomic unit minus
one makes the entire translation subgroup
commutators.
Every abelian quotient is therefore
cyclotomic.  Since the local cyclotomic
extension at three is totally ramified,
the unramified quotient is trivial.
It follows that \(L_n/\Q\) is totally
ramified at its unique three-prime.
As \(F/k\) is split there,
\[
                              F\cap L_n=k.
\tag{A.151.7}
\]
This proves that \(FL_n/L_n\) has degree
three and identifies its group with
\(\operatorname{Gal}(F/k)\) by restriction.

\emph{The exact local inclusions at both new support primes.}
Fix \(q\in\{11,59\}\), and write
\(M_q=3^s m\) with \(3\nmid m\).
The unit part of \(a\) in \(\Q_q\) has
compatible roots of every three-power
order in \(\Q_q\): the order \(q-1\)
of its residue-unit group is prime to
three, and the power maps on principal
units are invertible.
For \(n>s\), divide the chosen radical
\(a^{1/3^n}\), in a completion of \(L_n\),
by such a unit root.
The result, denoted \(\beta\), satisfies
\(\beta^{3^{n-s}}=\zeta q^m\) for
\(\zeta\in\mu_{3^s}\).
Choose \(\omega\in\mu_{3^n}\) with
\(\omega^{3^{n-s}}=\zeta^{-1}\), and
replace \(\beta\) by \(\omega\beta\).
Its \(3^{n-s}\)-th power is now literally
\(q^m\).
Choose integers \(A,B\) with
\(mA+3^{n-s}B=1\).
Raising that root to \(A\) and multiplying
by \(q^B\) gives
\(q^{1/3^{n-s}}\) in the same local field.
In particular
\[
                         q^{1/3}\in(L_n)_w
                         \quad(n>s, w\mid q).
\tag{A.151.8}
\]
The other rational support prime is a
unit in \(\Q_q\), hence is itself a
cube there.  Thus \(649^{1/3}\) lies
in every indicated completion of
\(L_n\) for \(n\geq n_0\).
This proves that \(FL_n/L_n\) splits
at both sets of support primes.
It splits at two and three by the
already proved local cubes in the
base fields.
At every other finite place,
\(F/k\) is unramified, and base
change preserves that assertion.
All archimedean places of \(L_n\)
are complex.
The claimed unramified and split
properties are now proved at every
place.

\emph{The actual Artin characters and their norm transitions.}
An unramified abelian extension split
at \(S_n\) is a quotient of the
\(S_n\)-ideal class group: its Artin
map kills principal ideals and every
prime inverted in the ring.
Restriction to \(F\) and (A.151.7)
therefore give the onto map
\(\chi_n\) in (A.151.2).
For a prime \(\mathfrak P\) of
\(L_{n+1}\) outside \(S_{n+1}\),
lying over \(\mathfrak p\) of
\(L_n\), ideal norm and Artin
restriction give
\[
 \begin{aligned}
 N\mathfrak P
       &=\mathfrak p^{f(\mathfrak P/\mathfrak p)},\\
 \chi_n(N\mathfrak P)
       &=\operatorname{Frob}_{\mathfrak p}(F L_n/L_n)
                         ^{f(\mathfrak P/\mathfrak p)}\big|_F
         =\chi_{n+1}(\mathfrak P).
 \end{aligned}
\tag{A.151.9}
\]
This is the norm--Artin diagram of
\cite[V, Proposition 3.3]{MilneCFT}, on
the literal common quotient \(F/k\).
Its abelian-extension hypothesis holds here:
\(L_{n+1}/L_n\) is the compositum of the
two cyclic cubic Kummer extensions obtained
by adjoining the next radical and the
next three-power roots of unity, and
\(FL_n/L_n\) is cyclic cubic.
Thus \(FL_{n+1}/L_n\) is abelian.
It descends to the class groups and
proves compatibility.

Fix any \(t\in\Z/3\).
Each finite fibre \(\chi_n^{-1}(t)\)
is nonempty, and its transition norms
land in the preceding fibre.
A member at any largest finite height
gives all lower members by norm.
The product of these finite compact
fibres therefore satisfies the finite
intersection property, proving the
surjectivity in (A.151.3).

For a finite-coefficient cohomological
description of the same character,
regard its unramified, locally split
class as an element of
\[
 \ker\!\left(H^1(G_{L_n,S_n},\Z/3)
                 \longrightarrow\prod_{w\in S_n}H^1((L_n)_w,\Z/3)\right).
\tag{A.151.10}
\]
It is nonzero by (A.151.7).
Shapiro maps it to the corresponding
localization kernel over \(k\) for
the actual permutation module
\(\operatorname{Ind}_{G_{L_n,S_n}}^{G_{k,S}}\Z/3\).
Finite Poitou--Tate duality
\cite[I, Theorem 4.10(a)]{BIR_MilneADT} gives
the dual nonzero degree-two
localization kernel for its Cartier
dual module.
This specifies a finite coefficient
map for the class-group phenomenon.
No nonzero image of a particular
one-variable target obstruction
under this pairing is asserted.

Finally, \(11\) and \(59\) are both
\(11\pmod{24}\).
Thus \(\gcd(24,q-1)=2\), and
\(-2\) is a square, hence a
twenty-fourth power, modulo each.
For \(M_1=M_2=24\), the rational
point \((X,Y)=(649^{-1},1)\) on
the old cover is a unit point at
every finite prime other than
eleven and fifty-nine.
At either of those primes, set
\(X=1\) and Hensel-lift a unit
twenty-fourth root of \(-1/2\).
The derivative is a unit because
neither prime divides twenty-four.
This proves the retained all-place
unit-point statement without changing
the class-group scope.
\end{proof}

The exact certificate
\path{code/two_inert_prime_class_limit_certificate.py}
and its output retain the ray quotient,
the split-three Hensel inequality,
the local root-exponent maps and the
actual old-cover residue witnesses.
